\subsection{DEFINITIONS OF THE PRODUCT OF A FAMILY OF SETS}

Let \((\mathbf{X}_{\iota})_{\iota \in \mathbf{I}}\) be a family of sets and let \(\mathbf{F}\) be a functional graph with domain \(\mathbf{I}\) such that \(\mathbf{F}(\iota) \in \mathbf{X}_{\iota}\) for each \(\iota \in \mathbf{I}\) . Then for each \(\iota \in \mathbf{I}\) we have \(\mathbf{F}(\iota) \in \mathbf{A} = \bigcup_{\iota \in \mathbf{I}} \mathbf{X}_{\iota}\) , and therefore \(\mathbf{F}\) is an element of \(\mathfrak{P}(\mathbf{I} \times \mathbf{A})\) . The functional graphs with the above property therefore form a subset of \(\mathfrak{P}(\mathbf{I} \times \mathbf{A})\) .

DEFINITION 1. Let \((\mathbf{X}_{\iota})_{\iota \in \mathbf{I}}\) be a family of sets. The set of functional graphs \(\mathbf{F}\) with domain I such that \(\mathrm{F}(\iota) \in \mathbf{X}_{\iota}\) for each \(\iota \in \mathbf{I}\) is called the product of the family of sets \((\mathbf{X}_{\iota})_{\iota \in \mathbf{I}}\) and is denoted by \(\prod_{\iota \in \mathbf{I}} \mathbf{X}_{\iota}\) . For each \(\iota \in \mathbf{I}\) , \(\mathbf{X}_{\iota}\) is called the factor of index \(\iota\) in the product \(\prod_{\iota \in \mathbf{I}} \mathbf{X}_{\iota}\) . The mapping \(\mathrm{F} \to \mathrm{F}(\iota)\) ( \(\mathrm{F} \in \prod_{\iota \in \mathbf{I}} \mathbf{X}_{\iota}\) , \(\mathrm{F}(\iota) \in \mathbf{X}_{\iota}\) ) is called the coordinate function (or projection) of index \(\iota\) , and is denoted by \(\operatorname{pr}_{\iota}\) .

\(F(t)\) is called the coordinate of index \(t\) (or projection of index \(t\) ) of F; the image \(pr_{t}\langle A\rangle\) of a subset A of \(\prod_{i\in I}X_{i}\) under the coordinate function of index \(t\) is called the projection of index \(t\) of A. It is easily verified that \(A\subset\prod_{i\in I}pr_{t}\langle A\rangle\) .

We shall often use the notation \((x_{i})_{i\in I}\) to denote an element of \(\prod_{i\in I}X_{i}\) (§3, no. 6).

If \(\mathbf{I} = \varnothing\) , the set \(\prod_{i\in \mathbf{I}}\mathbf{X}_i\) has only one element, namely the empty set (§3, no. 4, Example 1).

If all the factors \(X_{t}\) of the product \(\prod_{i\in I}X_{t}\) are equal to the same set E, we have \(\prod_{i\in I}X_{t}=E^{I}\) ; this follows immediately from the definitions.

If \((\mathbf{X}_t)_{t\in I}\) is an arbitrary family of sets and if \(\mathbf{E}\) is a set such that

\[
\bigcup_ {i \in I} X _ {i} \subset E,
\]

then Definition 1 shows that \(\prod_{i\in I}X_i\subset E^I\) ; there is therefore a one-to-one correspondence between \(\prod_{i\in I}X_i\) and a set of mappings of \(I\) into \(E\) (i.e., a subset of \(\mathcal{F}(I,E)\) ).

If \(\mathbf{I} = \{\alpha\}\) is a set consisting of a single element, we have

\[
\prod_ {i \in I} X _ {i} = X _ {\alpha} ^ {\{\alpha \}};
\]

the mapping \(\mathbf{F} \to \mathbf{F}(\alpha)\) is then a bijection (called canonical) of \(\prod_{i \in \{x\}} \mathbf{X}_i\) onto \(\mathbf{X}_\alpha\) .

Let A, B be sets and let \(\alpha, \beta\) be distinct objects (there exist two distinct objects, for example \(\phi\) and \(\{\phi\}\) ). Consider the graph \(\{(\alpha, A), (\beta, B)\}\) , which is clearly functional and is precisely the family \((X_{t})_{t \in [\alpha, \beta]}\) such that \(X_{\alpha} = A\) and \(X_{\beta} = B\) . For each pair \((x, y) \in A \times B\) , let \(f_{x,y}\) be the functional graph \(\{(\alpha, x), (\beta, y)\}\) . It is evident that the function \((x, y) \to f_{x,y}\) is a bijective mapping of \(A \times B\) onto \(\prod_{t \in [\alpha, \beta]} X_t\) ; the inverse of this bijection is \(g \to (g(\alpha), g(\beta))\) . These two bijections are called canonical. In what follows we shall use this one-to-one correspondence to deduce properties of the product of two sets from properties of the product of a family of sets.

Let \((\mathbf{X}_i)_{i\in I}\) be a family of sets each of which consists of a single element, say \(\mathbf{X}_i = \{a_i\}\) ; then the product \(\prod_{i\in I}\mathbf{X}_i\) is a set consisting of the single element \((a_i)_{i\in I}\) .

Let E be a set. The graphs of the constant mappings \(\iota \to x\) of I into E form a subset \(\Delta\) of the product \(\mathbf{E}^{\mathrm{I}}\) , called the diagonal. If \(\overline{x}\) denotes the graph of the mapping \(\iota \to x\) (where \(x \in \mathbf{E}\) ), the mapping \(x \to \overline{x}\) is a bijection of E onto \(\Delta\) , called the diagonal mapping.

PROPOSITION 4. Let \((\mathbf{X}_i)_{i\in I}\) be a family of sets, and let \(u\) be a bijection of a set \(K\) onto the index set \(I\) . If \(U\) is the graph of \(u\) , the mapping \(\mathbf{F} \to \mathbf{F} \circ U\) of \(\prod_{i\in I} X_i\) into \(\prod_{x\in K} X_{u(x)}\) is bijective.

Let

\[
\mathrm{A} = \bigcup_ {i \in \mathbf {I}} \mathrm{X} _ {i} = \bigcup_ {\boldsymbol {x} \in \mathbf {K}} \mathrm{X} _ {u (\boldsymbol {x})}
\]

(§4, Proposition 1). The mapping \(\mathbf{F} \to \mathbf{F} \circ \mathbf{U} (\mathbf{F} \in \mathbf{A}^{\mathbf{I}})\) is a bijection of \(\mathbf{A}^{\mathbf{I}}\) onto \(\mathbf{A}^{\mathbf{K}}\) (Proposition 2). It is evident that the condition ``for each \(\iota \in \mathbf{I}, \mathbf{F}(\iota) \in \mathbf{X}_{\iota}\) '' is equivalent to ``for each \(x \in K, (\mathbf{F} \circ \mathbf{U})(x) \in X_{u(x)}\) '', and the result follows.

